%=======================================================================
\section{Các Công việc Đã Hoàn thành}
\label{sec:completed}
%=======================================================================

\begin{enumerate}[label=\textbf{C\arabic*.},leftmargin=2.2em]
  \item \textbf{Thiết lập Pipeline River và Tích hợp Bộ Chỉ số Đánh giá} \hfill{\small(Giai đoạn 1 \& 2)}\\
  \begin{itemize}[leftmargin=1.2em,itemsep=2pt,topsep=1pt]
    \item Đã cài đặt River và xây dựng một khung thực thi pipeline thống nhất nhằm tối ưu hóa việc vận hành mô hình luồng, đánh giá và tích hợp thuật toán.
    \item Đã triển khai và đóng gói đầy đủ các chỉ số đánh giá phân cụm tương thích trực tiếp với các luồng dữ liệu của River:
    \begin{itemize}[leftmargin=1.2em,itemsep=1pt]
      \item \textbf{Chỉ số ngoại vi (External metrics):} Adjusted Rand Index (ARI), Normalized Mutual Information (NMI), Adjusted Mutual Information (AMI), Purity, Rand Index (RI).
      \item \textbf{Chỉ số nội tại (Internal metrics):} Hệ số Silhouette, Tổng bình phương sai số (SSQ/SSE), Chỉ số Davies--Bouldin (DB), Chỉ số Calinski--Harabasz (CH).
    \end{itemize}
  \end{itemize}

  \item \textbf{Triển khai Thuật toán Phân cụm Luồng và Chiến lược Ưu tiên} \hfill{\small(Giai đoạn 3)}\\
  \begin{itemize}[leftmargin=1.2em,itemsep=2pt,topsep=1pt]
    \item Đã viết mã nguồn cho 10 thuật toán phân cụm luồng. Các thuật toán này hiện chưa được kiểm chứng đối chiếu toàn diện với bài báo gốc; SVStream là thuật toán đầu tiên đang trong quá trình kiểm chứng (Mục~\ref{sec:svstream_reproduction}).
    \begin{itemize}[leftmargin=1.2em,itemsep=1pt]
      \item \textbf{Nhóm nền tảng, chuẩn mực và được trích dẫn rộng rãi:} DStream, HPStream, Birch~\cite{zhang1996birch}, RepStream~\cite{luhr2009repstream}, SVstream~\cite{wang2013svstream}. Birch vốn là phương pháp duyệt một lượt (one-pass) cho tập dữ liệu tĩnh kích thước lớn thay vì thuật toán phân cụm luồng thuần túy; thuật toán này được giữ lại làm baseline về cây CF mà các phương pháp luồng (như CluStream) kế thừa và phát triển.
      \item \textbf{Mô hình khám phá/khác:} DDStream, MrStream, DCUstream, improved\_Birch~\cite{ismael2014improvedbirch}, SOStream~\cite{isaksson2012sostream}.
    \end{itemize}
    \item \textbf{Chiến lược Triển khai:} Ưu tiên các thuật toán nền tảng vì chúng được trích dẫn nhiều, đặc tả thuật toán rõ ràng và có các thí nghiệm bài bản để đối chiếu mã nguồn. Mục~\ref{sec:svstream_reproduction} chỉ ra rằng ngay cả với nhóm này, việc tái lập chính xác 100\% cũng gặp nhiều hạn chế (dữ liệu không được công bố, đồ thị thiếu trục tọa độ, tham số mô tả mơ hồ); do đó, mỗi bản cài đặt cần được kiểm chứng đối chiếu kỹ lưỡng theo từng hình ảnh thí nghiệm trước khi được tính là hoàn tất.
  \end{itemize}

  \item \textbf{Nghiên cứu về Phương pháp luận Benchmark Siêu tham số cho Phân cụm Luồng} \hfill{\small(Phương pháp luận)}\\
  \begin{itemize}[leftmargin=1.2em,itemsep=2pt,topsep=1pt]
    \item \textbf{Thực trạng \& Thách thức:}
    \begin{itemize}[leftmargin=1.2em,itemsep=1pt]
      \item \emph{Vấn đề Quá khớp (Overfitting):} Hầu hết các tài liệu khoa học công bố các siêu tham số được tinh chỉnh thủ công riêng trên các tập dữ liệu do chính tác giả lựa chọn. Khi chạy benchmark trên các tập dữ liệu khác, hiệu năng suy giảm nghiêm trọng hoặc xuất hiện hiện tượng tăng vọt đột biến về thời gian chạy/bộ nhớ RAM. Trong học máy tổng quát, tối ưu hóa siêu tham số (HPO) đóng vai trò then chốt để xác định cấu hình hoạt động tối ưu và bền bỉ~\cite{franceschi2024hyperparameter}.
      \item \emph{So sánh Thiếu công bằng:} Việc lựa chọn tham số theo cảm tính (heuristic) hoặc tinh chỉnh trên chính tập dữ liệu kiểm thử tạo ra độ chệch đánh giá (evaluation bias), như đã được chỉ ra đối với các bộ phát hiện trôi dạt khái niệm trong~\cite{cerqueira2026framework}; chúng tôi nhận thấy điều tương tự cũng tồn tại trong phân cụm.
      \item \emph{Mục tiêu Cốt lõi:} Xác định cấu hình siêu tham số có khả năng hoạt động ổn định và khách quan trên nhiều tập dữ liệu đa dạng, thay vì tìm kiếm bộ tham số chỉ tối ưu cục bộ trên một tập dữ liệu đơn lẻ.
    \end{itemize}

    \item \textbf{Ý tưởng Tham khảo - Leave-One-Dataset-Out (LODO):} được đề xuất trong~\cite{cerqueira2026framework} dành cho benchmark các phương pháp phát hiện trôi dạt khái niệm, vốn chưa phải là quy trình chuẩn tắc cho bài toán phân cụm; chúng tôi đã vận dụng và điều chỉnh phương pháp này.
    \begin{itemize}[leftmargin=1.2em,itemsep=1pt]
      \item \emph{Cơ chế:} Sử dụng $n-1$ tập dữ liệu để tìm kiếm cấu hình tối ưu (thông qua tìm kiếm ngẫu nhiên - random search) và đánh giá trên tập dữ liệu giữ lại còn lại. Giúp quá trình tinh chỉnh hoàn toàn độc lập với dữ liệu kiểm thử.
      \item \emph{Ưu điểm:} Không phụ thuộc vào mô hình (model-agnostic), khảo sát được không gian tham số rộng, trực quan, đảm bảo tính tái lập.
      \item \emph{Nhược điểm:} Chi phí tính toán rất lớn, không gian tìm kiếm khổng lồ, phụ thuộc nhiều vào yếu tố may rủi hơn là tri thức miền chuyên sâu.
      \item \emph{Hạn chế đối với Phân cụm (quan sát của chúng tôi, không trích xuất từ tài liệu đã xuất bản):} LODO tiêu chuẩn gặp thất bại vì các tham số tĩnh tuyệt đối (ví dụ: $\epsilon$ cố định) không mang ý nghĩa hình học nhất quán giữa các tập dữ liệu có số chiều, thang đo và số lượng cụm khác biệt lớn (ví dụ: KDD99 có 23 cụm so với Shuttle chỉ có 2 cụm). Không thể tồn tại một giá trị tĩnh tuyệt đối nào phù hợp cho mọi trường hợp.
    \end{itemize}

    \item \textbf{Đề xuất của Chúng tôi - LODO có Chủ đích với Tham số Tương đối} (chưa đánh giá thực nghiệm):
    \begin{enumerate}[label=(\alph*),leftmargin=1.4em,itemsep=2pt]
      \item \textbf{Chỉ Tinh chỉnh Tham số Trọng yếu:} Tập trung duy nhất vào 1--2 tham số then chốt (hoặc nhiều hơn tùy thuật toán) ảnh hưởng trực tiếp đến cấu trúc cụm và độ phân giải:
      \begin{itemize}[leftmargin=1.2em,itemsep=1pt]
        \item \emph{Dựa trên tâm (KMeans, v.v.):} $k$ (số lượng cụm).
        \item \emph{Dựa trên mật độ (DenStream, v.v.):} $\epsilon$, \texttt{clustering\_threshold} (bán kính micro-cluster và tiêu chuẩn gộp/tách).
        \item \emph{Dựa trên lưới (DStream, v.v.):} \texttt{grid\_width} (ngưỡng mật độ ô lưới).
        \item \emph{Phân cấp (BIRCH, v.v.):} \texttt{threshold} (ngưỡng phân nhánh cây CF).
        \item \emph{Dựa trên đồ thị (RepStream, v.v.):} $\alpha$ (mật độ điểm đại diện).
      \end{itemize}
      \item \textbf{Sử dụng Tham số Tương đối qua Đặc trưng Thống kê Dữ liệu ($\text{param} = f(\text{stats})$):} Tinh chỉnh các hệ số co giãn tương đối thay vì các giá trị tuyệt đối tĩnh:
      \begin{itemize}[leftmargin=1.2em,itemsep=1pt]
        \item $\epsilon = c \cdot d$, với $d$ là khoảng cách trung vị từng cặp hoặc khoảng cách $k$-NN trung bình.
        \item $\text{grid\_width} = \text{range}(X) / g$, với $g$ là số lượng phân đoạn mong muốn trên mỗi chiều.
        \item $\text{threshold} = c \cdot \sigma$, với $\sigma$ là độ lệch chuẩn trung bình trên các chiều thuộc tính.
        \item $k$ và $\alpha$ gắn liền với số lượng lớp hoặc ước lượng mode phân bố trong mẫu.
      \end{itemize}
      \item \textbf{Lợi ích Cốt lõi:} Một hệ số co giãn duy nhất $c$ sẽ tự động thích ứng với mật độ dữ liệu (tạo $\epsilon$ lớn hơn cho dữ liệu thưa, nhỏ hơn cho dữ liệu dày), loại bỏ hoàn toàn nhược điểm của việc áp đặt các giá trị tuyệt đối cứng nhắc trên các tập dữ liệu không đồng nhất. Ý tưởng này chia sẻ tinh thần tương đồng với việc tinh chỉnh siêu tham số theo hướng dữ liệu trong phân cụm~\cite{fan2020hyperparameter}, nhưng bản thân phương pháp LODO với tham số tương đối là đề xuất riêng của chúng tôi và hiện chưa được đánh giá thực nghiệm.
    \end{enumerate}
  \end{itemize}
\end{enumerate}
